% ============================================================
% CAPITOLO 30 - PROGRESSIONI CON SETTIME
% ============================================================
\chapter{Progressioni con settime}
\label{chap:progressioni_settime}

Le settime aumentano la tensione direzionale delle progressioni. Questo capitolo riprende gli schemi a moduli dei Capitoli \ref{chap:progressioni_armoniche} e \ref{chap:moduli} e li realizza con le settime, poi tratta la catena di settime per quinte discendenti.

\section{Progressioni con settime}
\label{sec:progressioni_settime}

\begin{regola}[Progressioni di settime]
	\label{reg:progressioni_settime}
	Nelle progressioni le settime si possono impiegare in due modi:
	\begin{enumerate}
		\item \textbf{Settime alternate:} solo alcuni accordi della progressione hanno la settima (di solito il secondo accordo del modulo).
		\item \textbf{Settime consecutive:} ogni accordo è un accordo di settima.
	\end{enumerate}
	Valgono le regole di condotta delle progressioni a moduli (Regola \ref{reg:condotta_moduli}) con queste aggiunte:
	\begin{itemize}
		\item la settima è \textbf{preparata} (nota già presente nella stessa voce) e \textbf{scende di grado};
		\item la quinta diminuita scende di grado quando l'accordo seguente contiene la nota sotto;
		\item nell'accordo con la settima, quando la quinta è omessa, si raddoppia la fondamentale: gli accordi di una catena di settime alternano \textbf{completo e incompleto}.
	\end{itemize}
\end{regola}

Per ogni schema si danno tre versioni: solo triadi (già nel Capitolo \ref{chap:moduli}), solo settime, triadi e settime con la settima sul secondo accordo del modulo.

\begin{esempio}[I--IV discendente in Do maggiore: non modulante]
	\label{ex:prog_i4d_dia_sett}
	\textbf{Solo settime}
	\musica{es_cap17_i4d_dia_S}
	\textbf{Triadi e settime}
	\musica{es_cap17_i4d_dia_TS}
\end{esempio}

\begin{esempio}[I--IV discendente: modulante, sequenza reale di tre moduli]
	\label{ex:prog_i4d_rea_sett}
	\textbf{Solo settime}
	\musica{es_cap17_i4d_rea_S}
	\textbf{Triadi e settime}
	\musica{es_cap17_i4d_rea_TS}
\end{esempio}

\begin{esempio}[I--IV ascendente in Do maggiore: non modulante]
	\label{ex:prog_i4a_dia_sett}
	\textbf{Solo settime}
	\musica{es_cap17_i4a_dia_Sa}
	\musica{es_cap17_i4a_dia_Sb}
	\textbf{Triadi e settime}
	\musica{es_cap17_i4a_dia_TSa}
	\musica{es_cap17_i4a_dia_TSb}
\end{esempio}

\begin{esempio}[I--IV ascendente: modulante, sequenza reale di tre moduli]
	\label{ex:prog_i4a_rea_sett}
	\textbf{Solo settime}
	\musica{es_cap17_i4a_rea_S}
	\textbf{Triadi e settime}
	\musica{es_cap17_i4a_rea_TS}
\end{esempio}

\begin{esempio}[Progressione imitata I--IV discendente con settime]
	\label{ex:imitata_i4d_sett}
	\textbf{Solo settime}
	\musica{es_cap17_i4d_imi_S}
	\textbf{Triadi e settime}
	\musica{es_cap17_i4d_imi_TS}
\end{esempio}

\begin{osservazione}[Schemi non realizzabili con le settime]
	\label{oss:schemi_non_realizzabili}
	Gli schemi I--V, ascendente e discendente, e l'imitata I--IV ascendente non si realizzano con le settime. La settima dell'accordo dovrebbe scendere di grado, ma la nota sotto di essa non appartiene all'accordo seguente (nel \sigla{V}{7}{} il Fa dovrebbe scendere a Mi, assente nel II), oppure la ripetizione identica del modello produce parallele.
\end{osservazione}

\section{Catena di settime}
\label{subsec:catena_settime}

Una variante più sofisticata della progressione per quinte usa sistematicamente accordi di settima in stato fondamentale: la \textbf{catena di settime}. Il movimento delle voci è quasi automatico.

\begin{regola}[Progressione per quinte con accordi di settima (tutti 7)]
	\label{reg:progressione_7}
	Nella catena di settime:
	\begin{itemize}
		\item ogni accordo è una \textbf{settima} in stato fondamentale (numerica 7);
		\item il basso si muove per quinte discendenti (o quarte ascendenti);
		\item la settima di ogni accordo \textbf{è preparata} (compare come consonanza nell'accordo precedente) e \textbf{scende di grado} nell'accordo successivo;
		\item la settima di un accordo diventa la terza dell'accordo seguente, che sta una quinta sotto.
	\end{itemize}
\end{regola}

Lo schema delle voci è costante: la settima scende di grado, la quinta scende di grado sulla fondamentale seguente, la terza resta ferma e diventa settima. Gli accordi si alternano \textbf{completo e incompleto}. Nel basso il tritono Fa--Si è inevitabile nella scala diatonica.

\begin{esempio}[Catena di settime in Do maggiore e in La minore]
	\label{ex:catena_do_la}
	Do maggiore: \sigla{I}{7}{} -- \sigla{IV}{7}{} -- \sigla{VII}{7}{} -- \sigla{III}{7}{} -- \sigla{VI}{7}{} -- \sigla{II}{7}{} -- \sigla{V}{7}{} -- I.
	\musica{es_cap17_catena_do}
	La minore (naturale, con la sensibile alzata nella \sigla{V}{7}{}): \sigla{i}{7}{} -- \sigla{iv}{7}{} -- \sigla{VII}{7}{} -- \sigla{III}{7}{} -- \sigla{VI}{7}{} -- \sigla{iiø}{7}{} -- \sigla{V}{7}{} -- i.
	\musica{es_cap17_catena_la}
\end{esempio}
